\subsubsection{A finite Frobenius fiber excludes every bounded coupled lift}

The direct-sum hypothesis can be removed in the ordinary type-A
Soergel category. The extra input is an actual finite coinvariant
fiber and its algebraic duality. This argument concerns finite
Frobenius algebras; its use of that word has no connection to an
arithmetic Frobenius endomorphism.

\begin{finding}[\fid{F59}, bounded coupled-complex strengthening]
\label{f:boundedSoergelNoGo}
Let \(n\ge2\), let \(k\) have characteristic zero, and take the
standard permutation realization of \(S_n\times S_n\) on
\[
 A=k[x_1,\ldots,x_n,y_1,\ldots,y_n],
 \qquad \deg x_i=\deg y_i=2.
\]
For a simple adjacent transposition \(s\), no bounded complex of
ordinary Soergel bimodules has both Euler class
\[
 P_s=b_s\otimes b_s+(b_s-\delta)\otimes(b_s-\delta)
\]
and a homotopy isomorphism
\begin{equation}\label{eq:boundedSoergelquadratic}
 X\otimes_A X\simeq
 \bigl(A(2)\oplus A^{\oplus2}\oplus A(-2)\bigr)\otimes_A X.
\end{equation}
No direct-sum decomposition of \(X\) is assumed, and its bounded
differentials are unrestricted.
\end{finding}

\begin{proof}
\par\smallskip\noindent\textbf{A central finite-Frobenius lemma.}\quad
First let \(A\) be any connected nonnegatively graded polynomial
\(k\)-algebra and let \(I\subset A_+\) be homogeneous. Suppose
\(B=A/I\) is a finite-dimensional commutative local graded
Frobenius algebra: there is a homogeneous functional
\(\tau:B\to k\) with nondegenerate pairing \(\tau(ab)\).
Let \(X\) be a bounded graded bimodule complex with terms finite
free on both sides. Suppose \(I\) has homogeneous generators
which each act identically from the left and right on every term.
This implies \(IM=MI\) for each term \(M\), by moving the
generators across the bimodule; arbitrary elements of the ideal
need not themselves act centrally.

Assume the quadratic
\eqref{eq:boundedSoergelquadratic} and a nonzero Euler left
rank. We claim that \(X\) has cohomology only in degree zero,
with \(H^0X\) finite graded-free on the left.

Put \(Z=X/IX\) termwise. Because \(IX^j=X^jI\), its terms
are free on both sides over \(B\). Balanced tensor products give
\begin{equation}\label{eq:centralquotientmonoidal}
 (M/IM)\otimes_B(N/IN)
 \simeq(M\otimes_A N)/I(M\otimes_A N).
\end{equation}
Both sides quotient \(M\otimes_kN\) by the balancing relations
and by \(I\) in either factor; the equalities \(IM=MI\) and
\(IN=NI\) identify these last relations. The same formula holds
for total complexes. The quotient functor preserves chain
homotopies because it is additive. Therefore the actual
homotopy relation reduces to
\begin{equation}\label{eq:finiteFrobeniusquadratic}
 Z\otimes_B Z\simeq\mathscr D_B\otimes_B Z,
 \qquad\mathscr D_B=B(2)\oplus B^{\oplus2}\oplus B(-2).
\end{equation}
This step needs no flatness over an invariant subring and no
cohomology-and-base-change assertion. The free terms compute
the derived tensor products. Moreover,
\[
 \chi_k(Z)=\dim_kB\sum_j(-1)^j\operatorname{rank}_A X^j\ne0,
\]
so \(Z\) has nonzero cohomology.

\par\smallskip\noindent\textbf{Both cohomological extremes.}\quad
For nonzero finite right and left modules \(M,N\) over the
commutative local algebra \(B\),
\(M\otimes_BN\ne0\): modulo its maximal ideal the tensor
product is the tensor product of two nonzero residue vector
spaces, by Nakayama. The right action on \(M\) and left action
on \(N\) are the ones used here; their other bimodule actions
need not agree.

Let \(r\) be the largest degree with \(H^rZ\ne0\).
Truncation of bounded-above complexes gives
\begin{equation}\label{eq:FrobeniusTopTensor}
 H^{2r}(Z\otimes_B^L Z)
       \simeq H^rZ\otimes_B H^rZ\ne0.
\end{equation}
Indeed every lower truncation contributes only below degree
\(2r\), and tensoring the two top modules has ordinary tensor
product in that degree and Tor only below it. No higher
cohomology occurs. The right side of
\eqref{eq:finiteFrobeniusquadratic} has largest degree \(r\),
so \(2r=r\) and \(r=0\).

For the other extreme use the graded bimodule dual
\[
 Z^\vee=\underline{\operatorname{Hom}}_{-B}(Z,B),
\]
where the underline includes all internal degrees. The
Frobenius functional identifies this functor with \(k\)-linear
duality, up to a fixed internal shift. Explicitly, a right-linear
map \(\phi\) gives \(\tau\circ\phi\); its inverse is
characterized by
\[
                  \tau(\phi(m)b)=\psi(mb)\qquad(b\in B).
\]
The nondegenerate form supplies existence and uniqueness.
The two bimodule actions interchange under this identification.
Hence the functor is exact, detects nonzero finite modules,
preserves freeness on both sides, and gives
\[
 H^i(Z^\vee)=
 \underline{\operatorname{Hom}}_{-B}(H^{-i}Z,B).
\]
Finite dual bases give the tensor-dual isomorphism, with order
reversed,
\[
             (Z\otimes_B Z)^\vee
                 \simeq Z^\vee\otimes_B Z^\vee.
\]
For complexes this uses the usual dual differential and Koszul
signs. It is bimodule duality, not a symmetry assertion.
The coefficient object \(\mathscr D_B\) is self-dual: its
shifts two and minus two exchange. It commutes with every object
as a direct sum of shifted units. Dualizing
\eqref{eq:finiteFrobeniusquadratic} thus gives the same
quadratic for \(Z^\vee\). The upper-bound argument applied to
this nonzero dual forces its greatest cohomology degree to zero,
so the least cohomology degree of \(Z\) is zero as well.
Thus \(Z\) has only \(H^0Z\).

\par\smallskip\noindent\textbf{Fiber projectivity and minimality.}\quad
The module \(H^0Z\) has finite projective dimension on the
left. To see this, split off the positive acyclic tail of the
bounded free complex from the top; its remaining negative tail
is a finite projective resolution. A Frobenius algebra is
self-injective, since its regular module is, up to grading,
the \(k\)-linear dual of itself and
\(\operatorname{Hom}_B(-,D_kB)=D_k(-)\) is exact.
Every projective term is therefore injective. A finite
projective resolution splits from its left end, and repeated
splitting proves that \(H^0Z\) is projective. These are
left-module splittings; no bimodule contractibility is inferred.

Forget the right action on \(X\). Repeatedly cancel any nonzero
constant entry of its graded left-free differential matrices.
Such an entry is a homogeneous unit, and graded row and column
operations split off a contractible pair. Each cancellation
removes two generators, so the process terminates with a bounded
minimal left-free complex \(F\) homotopy equivalent to \(X\).
Its differential entries all lie in \(A_+\).

The left-module homotopy equivalence descends under quotient.
Thus \(F/IF\) is a minimal left \(B\)-free complex homotopy
equivalent to \(Z\), and hence to its projective \(H^0\) in
degree zero. Here minimality holds because the images of the
differential entries lie in \(B_+\), the maximal ideal of the
finite connected graded algebra. Tensor this homotopy
equivalence with \(B/B_+\). The differential on the left is
zero, and its degree-\(j\) vector space has dimension
\(\operatorname{rank}_A F^j\). A zero-differential complex
homotopy equivalent to a degree-zero vector space has no terms
outside degree zero. Hence \(F^j=0\) there. This proves that
\(X\) has only \(H^0\), finite graded-free on the left, as
claimed by the central finite-Frobenius lemma.

\par\smallskip\noindent\textbf{Verifying the type-A fiber.}\quad
For one permutation factor put
\[
 R_n=k[x_1,\ldots,x_n],\quad
 I_n=(e_1,\ldots,e_n),\quad B_n=R_n/I_n,
\]
where \(e_j\) is elementary symmetric of ordinary degree \(j\).
Each \(e_j\) lies in \(R_n^s\) for every simple reflection,
so it acts identically from the two sides on
\(R_n\otimes_{R_n^s}R_n(1)\). This property passes to
tensor products, shifts and direct summands.

The finite Frobenius property is a standard coinvariant fact;
the following verification makes its hypotheses explicit.
Let
\[
 g_j=h_j(x_j,\ldots,x_n),\qquad J_n=(g_1,\ldots,g_n),
\]
where \(h_j\) is the complete homogeneous polynomial of
ordinary degree \(j\) in the displayed tail. We first prove
\(I_n=J_n\). Modulo \(I_n\),
\[
 \prod_{i=1}^n(1-x_it)=1,
 \qquad
 \prod_{i=j}^n(1-x_it)^{-1}=\prod_{i<j}(1-x_it).
\]
The right side has degree at most \(j-1\), so the coefficient
in degree \(j\) gives \(g_j=0\). Thus \(J_n\subset I_n\).
Conversely work modulo \(J_n\) and put
\(H_j(t)=\prod_{i=j}^n(1-x_it)^{-1}\).
Descending induction shows that \(H_j\) is a polynomial of
degree at most \(j-1\). The recurrence
\(H_j=H_{j+1}/(1-x_jt)\), the vanishing coefficient
\(g_j=0\), and the induction hypothesis first kill the
degree-\(j\) coefficient and then every higher coefficient.
The induction starts with the empty-tail series one.
It ends with \(H_1=1\), so all elementary symmetric
coefficients of its inverse vanish. Hence \(I_n\subset J_n\).

Each \(g_j\) is monic of degree \(j\) in \(x_j\), with
coefficients in the later variables. Adjoin the variables in
the order \(x_n,x_{n-1},\ldots,x_1\), imposing
\(g_n,g_{n-1},\ldots,g_1\) successively.
Division by a monic polynomial works over any coefficient
ring, so the resulting algebra has basis
\[
 \left\{\prod_{j=1}^n x_j^{a_j}:0\le a_j<j\right\},
 \qquad \dim_k B_n=n!.
\]
It is finite connected graded; its positive-degree ideal is
nilpotent and is its unique maximal ideal.

If \(T\) is a finite commutative Frobenius algebra and
\(C=T[z]/(z^j+a_1z^{j-1}+\cdots+a_j)\), extraction of
the coefficient of \(z^{j-1}\) makes \(C\) a Frobenius
extension of \(T\). In the basis \(1,z,\ldots,z^{j-1}\),
its multiplication pairing is zero below the reverse
diagonal and has ones on that diagonal. Its determinant is
therefore plus or minus one. Composing this perfect
\(T\)-pairing with the Frobenius functional on \(T\)
gives a perfect \(k\)-pairing. Induction through the monic
extensions proves that \(B_n\) is Frobenius. The relations
are homogeneous, so the resulting functional is homogeneous,
supported on the top basis monomial.

For the two factors use
\[
 I=I_n\otimes R_n+R_n\otimes I_n\subset A,
 \qquad B=B_n\otimes_k B_n.
\]
The tensor product of the two functionals is nondegenerate.
This is a finite local graded Frobenius algebra of dimension
\((n!)^2\), and the displayed generators of \(I\) act
centrally on every Soergel term. Every such term is finite
graded-free on both sides: Bott--Samelson terms are, and a
graded projective summand is graded-free by lifting a
homogeneous residue basis and applying graded Nakayama.
This verifies the central lemma's fiber and term hypotheses
in every type-A rank. For \(n=2\), after eliminating the
two invariant coordinates, the fiber is simply
\(k[x,y]/(x^2,y^2)\), with functional extracting \(xy\).

\par\smallskip\noindent\textbf{Applying the lemma to the actual class.}\quad
The generic-amplitude argument in the preceding finding
uses only the actual class and the quadratic. It gives
\begin{equation}\label{eq:boundedSoergelgenericsupport}
 H^i(X_Q)=0\quad(i\ne0),\qquad
 H^0(X_Q)\simeq Q_e^{\oplus2}\oplus Q_\tau^{\oplus2},
 \qquad\tau=(s,s).
\end{equation}
In particular the Euler left rank is four. The reduced
complex has Euler dimension \(4\dim_kB\ne0\), so the
central finite-Frobenius lemma applies. Consequently
\(X\) has only \(H^0=M\), and \(M\) is finite
graded-free on the left.

\par\smallskip\noindent\textbf{The two anti-invariant coordinates.}\quad
Change coordinates to write
\[
 A=k[x,y,z_1,\ldots,z_m],\qquad
 \tau(x)=-x,\quad\tau(y)=-y,\quad\tau(z_j)=z_j.
\]
Here \(x,y\) are the two independent simple roots. In a
homogeneous left basis of \(M\), let \(U,V\) be right
multiplication by \(x,y\). Generic support
\eqref{eq:boundedSoergelgenericsupport} gives
\[
                       yU=xV,\qquad U^2=x^2I.
\]
Left freeness descends these matrix identities to \(A\).
Since \(x,y\) are relatively prime, \(x\) divides every
entry of \(U\), so \(U=xZ\), \(V=yZ\), and \(Z^2=I\).
The operator \(Z\) has degree zero. Each fixed coordinate
acts on the right by its left scalar, by generic equality
and left torsion-freeness. Thus
\((I+Z)/2,(I-Z)/2\) are degree-zero bimodule idempotents.

Their images are graded projective left modules and hence
graded-free. Explicitly, lift a homogeneous residue basis,
apply graded Nakayama for a free surjection, split it in the
graded category, and apply graded Nakayama again to its
complementary kernel with zero residue quotient. On these
two images the right actions are respectively identity and
\(\tau\). Therefore
\begin{equation}\label{eq:boundedSoergelgraphsplit}
 M\simeq A_e^{\oplus a(v)}\oplus A_\tau^{\oplus b(v)},
 \qquad a(v),b(v)\in\mathbb N[v,v^{-1}],\quad b(v)\ne0.
\end{equation}
The last condition follows from the generic multiplicity
two of the \(\tau\) graph. This splitting also proves
right freeness.

The ordinary derived bimodule truncation identifies \(X\)
with \(M\). Since the latter is free on both sides, its
derived tensor square is its ordinary tensor square.
Thus \eqref{eq:boundedSoergelquadratic} induces a module
quadratic for \(M\). This conclusion uses derived
truncation; it does not misinterpret the earlier
left-module homotopy splitting as a bimodule map.

There are no maps between the identity and \(\tau\)
graph lines: any such scalar is annihilated by \(2x\)
in the polynomial domain. Comparing the two graded
multiplicities in the module quadratic gives
\[
 a(v)^2+b(v)^2=\delta^2a(v),\qquad
 2a(v)b(v)=\delta^2b(v).
\]
The Laurent ring is a domain and \(b(v)\ne0\), so
\[
                       2a(v)=v^2+2+v^{-2}.
\]
The outer coefficients on the right are odd, whereas
every coefficient on the left is even. This contradiction
proves the finding.
\end{proof}

The result closes the bounded coupled-complex repair in
the ordinary type-A Soergel category with the original
ambient class and the cohomologically degree-zero
coefficient for \(\delta^2\). A higher nonstandard
relation cannot repair failure of this quadratic.
The proof does not cover an unbounded completion, a
coefficient complex with the same Euler class but other
cohomological degrees, a localization changing the finite
fiber, or another monoidal category and Grothendieck map.
In particular the finite-state kernels of \fid{F56} are
outside these hypotheses. They retain the specified
tensor action without being bounded ordinary-Hecke
lifts of this ambient class. No arithmetic trace law,
source positivity conjecture or VIII basis identification
is used in the obstruction.

% Separate paragraph for Part III.
\paragraph{The finite-fiber strengthening in Grothendieck's language.}
The affine quotient and bimodule sheaves use EGA~I,
Theorem~1.3.7, p.~86, Corollary~1.3.8, p.~87, and
Theorem~1.4.1, pp.~90--92 \cite{EGA}.
The relevant perfect complexes have finite free terms
on the displayed tensor sides, in the framework of
SGA~6, I, Definition~4.7, pp.~121--122, and
Example~4.8, p.~122 \cite{SGA6}.
Preservation of the relation under the invariant-origin
quotient is proved by
\eqref{eq:centralquotientmonoidal}; it is not a generic
cohomology-and-base-change assertion. Exact duality of
the finite Frobenius algebra supplies the lower
cohomological bound. This algebraic duality is distinct
from arithmetic Frobenius and Weil purity.
The actual graded multiplicities in
\eqref{eq:boundedSoergelgraphsplit} use the effective
torus-character language of SGA~3, I, 4.7.2--3 and
cote~90, scan~17, and the effective/virtual distinction
of SGA~6, V, \S\S2--3 and cote~133, scan~18
\cite{SGA3,SGA6,Archive}.
No \(\lambda\)-ring structure on a noncommutative
Grothendieck algebra is asserted. The contradiction
uses a proved actual module splitting, not merely
a virtual Euler identity. The standard type-A
coinvariant fact was verified explicitly above; the
new categorical obstruction is not attributed to
any archival manuscript.
